\section{AMC 10 10A 2023}
        \begin{enumerate}
        	\item (\textit{AMC 10 10A 2023}) Cities $A$ and $B$ are $45$ miles apart. Alicia lives in $A$ and Beth lives in $B$. Alicia bikes towards $B$ at $18$ miles per hour. Leaving at the same time, Beth bikes toward $A$ at $12$ miles per hour. How many miles from City $A$ will they be when they meet?


 $\textbf{(A) }20\qquad\textbf{(B) }24\qquad\textbf{(C) }25\qquad\textbf{(D) }26\qquad\textbf{(E) }27$


	\item (\textit{AMC 10 10A 2023}) The weight of $\frac{1}{3}$ of a large pizza together with $3 \frac{1}{2}$ cups of orange slices is the same weight of $\frac{3}{4}$ of a large pizza together with $\frac{1}{2}$ cups of orange slices. A cup of orange slices weigh $\frac{1}{4}$ of a pound. What is the weight, in pounds, of a large pizza?


 $\textbf{(A) }1\frac{4}{5}\qquad\textbf{(B) }2\qquad\textbf{(C) }2\frac{2}{5}\qquad\textbf{(D) }3\qquad\textbf{(E) }3\frac{3}{5}$


	\item (\textit{AMC 10 10A 2023}) How many positive perfect squares less than $2023$ are divisible by $5$? 


 $\textbf{(A) }8\qquad\textbf{(B) }9\qquad\textbf{(C) }10\qquad\textbf{(D) }11\qquad\textbf{(E) }12$


	\item (\textit{AMC 10 10A 2023}) A quadrilateral has all integer side lengths, a perimeter of $26$, and one side of length $4$. What is the greatest possible length of one side of this quadrilateral?


 $\textbf{(A) }9\qquad\textbf{(B) }10\qquad\textbf{(C) }11\qquad\textbf{(D) }12\qquad\textbf{(E) }13$


	\item (\textit{AMC 10 10A 2023}) How many digits are in the base-ten representation of $8^5 \cdot 5^{10} \cdot 15^5$?


 $\textbf{(A) } 14 \qquad\textbf{(B) }15 \qquad\textbf{(C) }16 \qquad\textbf{(D) }17 \qquad\textbf{(E) } 18$


	\item (\textit{AMC 10 10A 2023}) An integer is assigned to each vertex of a cube. The value of an edge is defined to be the sum of the values of the two vertices it touches, and the value of a face is defined to be the sum of the values of the four edges surrounding it. The value of the cube is defined as the sum of the values of its six faces. Suppose the sum of the integers assigned to the vertices is $21$. What is the value of the cube? 


 $\textbf{(A) } 42 \qquad \textbf{(B) } 63 \qquad \textbf{(C) } 84 \qquad \textbf{(D) } 126 \qquad \textbf{(E) } 252$


	\item (\textit{AMC 10 10A 2023}) Janet rolls a standard $6$-sided die $4$ times and keeps a running total of the numbers she rolls. What is the probability that at some point, her running total will equal $3?$


 $\textbf{(A) }\frac{2}{9}\qquad\jtextbf{(B) }\frac{49}{216}\qquad\textbf{(C) }\frac{25}{108}\qquad\textbf{(D) }\frac{17}{72}\qquad\textbf{(E) }\frac{13}{54}$ (Error compiling LaTeX. Unknown error\_msg)


 

	\item (\textit{AMC 10 10A 2023}) Barb the baker has developed a new temperature scale for her bakery called the Breadus scale, which is a linear function of the Fahrenheit scale. Bread rises at $110$ degrees Fahrenheit, which is $0$ degrees on the Breadus scale. Bread is baked at $350$ degrees Fahrenheit, which is $100$ degrees on the Breadus scale. Bread is done when its internal temperature is $200$ degrees Fahrenheit. What is this, in degrees, on the Breadus scale?


 $\textbf{(A) }33\qquad\textbf{(B) }34.5\qquad\textbf{(C) }36\qquad\textbf{(D) }37.5\qquad\textbf{(E) }39$


	\item (\textit{AMC 10 10A 2023}) A digital display shows the current date as an $8$-digit integer consisting of a $4$-digit year, followed by a $2$-digit month, followed by a $2$-digit date within the month. For example, Arbor Day this year is displayed as $20230428.$ For how many dates in $2023$ does each digit appear an even number of times in the $8$-digital display for that date?


 $\textbf{(A)}~5\qquad\textbf{(B)}~6\qquad\textbf{(C)}~7\qquad\textbf{(D)}~8\qquad\textbf{(E)}~9$


	\item (\textit{AMC 10 10A 2023}) Maureen is keeping track of the mean of her quiz scores this semester. If Maureen scores an $11$ on the next quiz, her mean will increase by $1$. If she scores an $11$ on each of the next three quizzes, her mean will increase by $2$. What is the mean of her quiz scores currently? 


 $\textbf{(A) }4\qquad\textbf{(B) }5\qquad\textbf{(C) }6\qquad\textbf{(D) }7\qquad\textbf{(E) }8$


	\item (\textit{AMC 10 10A 2023}) A square of area $2$ is inscribed in a square of area $3$, creating four congruent triangles, as shown below. What is the ratio of the shorter leg to the longer leg in the shaded right triangle? 


 
                    \begin{figure}[H]
                        \centering
                        \includegraphics[width=0.5\textwidth]{Problem/AMC_10_10A_2023_11_Problem_Figure_1.png}
                    \end{figure}
                    
$\textbf{(A) }\frac15\qquad\textbf{(B) }\frac14\qquad\textbf{(C) }2-\sqrt3\qquad\textbf{(D) }\sqrt3-\sqrt2\qquad\textbf{(E) }\sqrt2-1$


	\item (\textit{AMC 10 10A 2023}) How many three-digit positive integers $N$ satisfy the following properties?


 
            \begin{itemize}
        \item The number $N$ is divisible by $7$.
            \end{itemize}
            
            \begin{itemize}
        \item The number formed by reversing the digits of $N$ is divisible by $5$.
            \end{itemize}
            $\textbf{(A) } 13 \qquad \textbf{(B) } 14 \qquad \textbf{(C) } 15 \qquad \textbf{(D) } 16 \qquad \textbf{(E) } 17$


	\item (\textit{AMC 10 10A 2023}) Abdul and Chiang are standing $48$ feet apart in a field. Bharat is standing in the same field as far from Abdul as possible so that the angle formed by his lines of sight to Abdul and Chiang measures $60^\circ$. What is the square of the distance (in feet) between Abdul and Bharat?


 $\textbf{(A) } 1728 \qquad \textbf{(B) } 2601 \qquad \textbf{(C) } 3072 \qquad \textbf{(D) } 4608 \qquad \textbf{(E) } 6912$


	\item (\textit{AMC 10 10A 2023}) A number is chosen at random from among the first $100$ positive integers, and a positive integer divisor of that number is then chosen at random. What is the probability that the chosen divisor is divisible by $11$?


 $\textbf{(A)}~\frac{4}{100}\qquad\textbf{(B)}~\frac{9}{200} \qquad \textbf{(C)}~\frac{1}{20} \qquad\textbf{(D)}~\frac{11}{200}\qquad\textbf{(E)}~\frac{3}{50}$


	\item (\textit{AMC 10 10A 2023}) An even number of circles are nested, starting with a radius of $1$ and increasing by $1$ each time, all sharing a common point. The region between every other circle is shaded, starting with the region inside the circle of radius $2$ but outside the circle of radius $1.$ An example showing $8$ circles is displayed below. What is the least number of circles needed to make the total shaded area at least $2023\pi$? 


 
                    \begin{figure}[H]
                        \centering
                        \includegraphics[width=0.5\textwidth]{Problem/AMC_10_10A_2023_15_Problem_Figure_1.png}
                    \end{figure}
                    
$\textbf{(A) } 46 \qquad \textbf{(B) } 48 \qquad \textbf{(C) } 56 \qquad \textbf{(D) } 60 \qquad \textbf{(E) } 64$


	\item (\textit{AMC 10 10A 2023}) In a table tennis tournament every participant played every other participant exactly once. Although there were twice as many right-handed players as left-handed players, the number of games won by left-handed players was $40\%$ more than the number of games won by right-handed players. (There were no ties and no ambidextrous players.) What is the total number of games played?


 $\textbf{(A) }15\qquad\textbf{(B) }36\qquad\textbf{(C) }45\qquad\textbf{(D) }48\qquad\textbf{(E) }66$


	\item (\textit{AMC 10 10A 2023}) Let $ABCD$ be a rectangle with $AB = 30$ and $BC = 28$. Point $P$ and $Q$ lie on $\overline{BC}$ and $\overline{CD}$ respectively so that all sides of $\triangle{ABP}, \triangle{PCQ},$ and $\triangle{QDA}$ have integer lengths. What is the perimeter of $\triangle{APQ}$?


 $\textbf{(A) } 84 \qquad \textbf{(B) } 86 \qquad \textbf{(C) } 88 \qquad \textbf{(D) } 90 \qquad \textbf{(E) } 92$


	\item (\textit{AMC 10 10A 2023}) A rhombic dodecahedron is a solid with $12$ congruent rhombus faces. At every vertex, $3$ or $4$ edges meet, depending on the vertex. How many vertices have exactly $3$ edges meet?


 $\textbf{(A) }5\qquad\textbf{(B) }6\qquad\textbf{(C) }7\qquad\textbf{(D) }8\qquad\textbf{(E) }9$


	\item (\textit{AMC 10 10A 2023}) The line segment formed by $A(1, 2)$ and $B(3, 3)$ is rotated to the line segment formed by $A'(3, 1)$ and $B'(4, 3)$ about the point $P(r, s)$. What is $|r-s|$?


 $\textbf{(A) } \frac{1}{4} \qquad \textbf{(B) } \frac{1}{2} \qquad \textbf{(C) } \frac{3}{4}   \qquad \textbf{(D) } \frac{2}{3} \qquad   \textbf{(E) } 1$


	\item (\textit{AMC 10 10A 2023}) Each square in a $3\times3$ grid of squares is colored red, white, blue, or green so that every $2\times2$ square contains one square of each color. One such coloring is shown on the right below. How many different colorings are possible? 


 
                    \begin{figure}[H]
                        \centering
                        \includegraphics[width=0.5\textwidth]{Problem/AMC_10_10A_2023_20_Problem_Figure_1.png}
                    \end{figure}
                    
$\textbf{(A) }24\qquad\textbf{(B) }48\qquad\textbf{(C) }60\qquad\textbf{(D) }72\qquad\textbf{(E) }96$


	\item (\textit{AMC 10 10A 2023}) Let $P(x)$ be the unique polynomial of minimal degree with the following properties:


 
            \begin{itemize}
        \item $P(x)$ has a leading coefficient $1$,
            \end{itemize}
            
            \begin{itemize}
        \item $1$ is a root of $P(x)-1$,
            \end{itemize}
            
            \begin{itemize}
        \item $2$ is a root of $P(x-2)$,
            \end{itemize}
            
            \begin{itemize}
        \item $3$ is a root of $P(3x)$, and
            \end{itemize}
            
            \begin{itemize}
        \item $4$ is a root of $4P(x)$.
            \end{itemize}
            The roots of $P(x)$ are integers, with one exception. The root that is not an integer can be written as $\frac{m}{n}$, where $m$ and $n$ are relatively prime integers. What is $m+n$?


 $\textbf{(A) }41\qquad\textbf{(B) }43\qquad\textbf{(C) }45\qquad\textbf{(D) }47\qquad\textbf{(E) }49$


	\item (\textit{AMC 10 10A 2023}) Circle $C_1$ and $C_2$ each have radius $1$, and the distance between their centers is $\frac{1}{2}$. Circle $C_3$ is the largest circle internally tangent to both $C_1$ and $C_2$. Circle $C_4$ is internally tangent to both $C_1$ and $C_2$ and externally tangent to $C_3$. What is the radius of $C_4$? 


 
                    \begin{figure}[H]
                        \centering
                        \includegraphics[width=0.5\textwidth]{Problem/AMC_10_10A_2023_22_Problem_Figure_1.png}
                    \end{figure}
                    
$\textbf{(A) } \frac{1}{14} \qquad \textbf{(B) } \frac{1}{12} \qquad \textbf{(C) } \frac{1}{10} \qquad \textbf{(D) } \frac{3}{28} \qquad \textbf{(E) } \frac{1}{9}$


	\item (\textit{AMC 10 10A 2023}) If the positive integer $c$ has positive integer divisors $a$ and $b$ with $c = ab$, then $a$ and $b$ are said to be $\textit{complementary}$ divisors of $c$. Suppose that $N$ is a positive integer that has one complementary pair of divisors that differ by $20$ and another pair of complementary divisors that differ by $23$. What is the sum of the digits of $N$?


 $\textbf{(A) } 9 \qquad \textbf{(B) } 13\qquad \textbf{(C) } 15 \qquad \textbf{(D) } 17 \qquad \textbf{(E) } 19$


	\item (\textit{AMC 10 10A 2023}) Six regular hexagonal blocks of side length $1$ unit are arranged inside a regular hexagonal frame. Each block lies along an inside edge of the frame and is aligned with two other blocks, as shown in the figure below. The distance from any corner of the frame to the nearest vertex of a block is $\frac{3}{7}$ unit. What is the area of the region inside the frame not occupied by the blocks?


 
                    \begin{figure}[H]
                        \centering
                        \includegraphics[width=0.5\textwidth]{Problem/AMC_10_10A_2023_24_Problem_Figure_1.png}
                    \end{figure}
                    
$\textbf{(A)}~\frac{13 \sqrt{3}}{3}\qquad\textbf{(B)}~\frac{216 \sqrt{3}}{49}\qquad\textbf{(C)}~\frac{9 \sqrt{3}}{2} \qquad\textbf{(D)}~ \frac{14 \sqrt{3}}{3}\qquad\textbf{(E)}~\frac{243 \sqrt{3}}{49}$


	\item (\textit{AMC 10 10A 2023}) If $A$ and $B$ are vertices of a polyhedron, define the distance $d(A, B)$ to be the minimum number of edges of the polyhedron one must traverse in order to connect $A$ and $B$. For example, $\overline{AB}$ is an edge of the polyhedron, then $d(A, B) = 1$, but if $\overline{AC}$ and $\overline{CB}$ are edges and $\overline{AB}$ is not an edge, then $d(A, B) = 2$. Let $Q$, $R$, and $S$ be randomly chosen distinct vertices of a regular icosahedron (regular polyhedron made up of $20$ equilateral triangles). What is the probability that $d(Q, R) > d(R, S)$?


 $\textbf{(A) }\frac{7}{22}\qquad\textbf{(B) }\frac{1}{3}\qquad\textbf{(C) }\frac{3}{8}\qquad\textbf{(D) }\frac{5}{12}\qquad\textbf{(E) }\frac{1}{2}$


\end{enumerate}